\chapter{Metodologia}

O presente capítulo descreve a metodologia empregada no desenvolvimento deste projeto, que se caracteriza pela união de diferentes áreas técnicas, como
a telecomunicação, sistemas embarcados, software/computação e tráfego urbano. A metodologia adotada é baseada em uma abordagem de pesquisa aplicada,
com foco na implementação prática por meio de simulações e validação dos conceitos estudados.

Um ponto crucial para o desenvolvimento deste trabalho é o uso de testes em ambiente virtual, que nos permite prever e analisar o comportamento do
sistema e encontrar possíveis falhas ou melhorias antes de uma implementação física. Validar virtualmente sistemas complexos é fundamental para
diminuir custos, otimizar recursos e limitar riscos, permitindo estudar cenários variados que podem ser encontrados no mundo real.

Para colocar esse estudo em prática, utiliza-se uma combinação de ferramentas de simulação, sendo o SUMO para modelagem de tráfego, o OMNeT++ para
simulação de redes e o Veins para integração entre ambos. A metodologia inclui a coleta de dados, modelagem do ambiente de tráfego, implementação
dos protocolos de comunicação V2X, execução das simulações e análise dos resultados obtidos.

\section{Análise do ambiente de tráfego}

\begin{alineas}
    \item Caracterização do ambiente de tráfego urbano;
    \item Requisitos de conectividade;
\end{alineas}

\section{Especificações da arquitetura C-V2X}
\begin{alineas}
    \item Padrões de comunicação;
    \item Mensagens cooperativas;
\end{alineas}

\section{Ambiente de simulação}
\begin{alineas}
    \item Configuração do SUMO;
    \item Configuração do OMNeT++;
    \item Integração com Veins.
\end{alineas}

\section{Procedimentos de validação}
\begin{alineas}
    \item explicar que a analise de impactos consiste em uma comparação entre cenários com e sem V2X;
\end{alineas}